% --- Archon physics formalization source begin ---
% archon:physics
% archon:covers PhyXMiniProblems/problem_phyx_mini_0532.lean
% archon:source-report reports/phyx_mini/problem_phyx_mini_0532.source.json

\chapter{Physics problem phyx\_mini\_0532}
\label{ch:PhyXMiniProblems_problem_phyx_mini_0532}

\paragraph{Problem source.}
\#\# Physical scenario

In the Davisson-Germer experiment, 54.0-eV electrons were diffracted from a nickel lattice. If the first maximum in the diffraction pattern was observed at \textbackslash{}( \textbackslash{}phi = 50.0\textasciicircum{}\textbackslash{}circ \textbackslash{}) as shown in figure.

\#\# Question

What was the lattice spacing \textbackslash{}( a \textbackslash{}) between the vertical columns of atoms in the figure?

\#\# Answer choices

- A: A: \textbackslash{}( 0.306 \textbackslash{}text\{ nm\} \textbackslash{})

- B: B: \textbackslash{}( 0.359 \textbackslash{}text\{ nm\} \textbackslash{})

- C: C: \textbackslash{}( 0.118 \textbackslash{}text\{ nm\} \textbackslash{})

- D: D: \textbackslash{}( 0.218 \textbackslash{}text\{ nm\} \textbackslash{})

\#\# Figure caption (auxiliary; use the image as primary evidence)

The image illustrates an electron diffraction setup. It shows a regular lattice of spherical atoms arranged in rows and columns on a plane. 

- **Objects and Arrangement:**
  - Spherical atoms are depicted as evenly spaced circles, arranged in a regular geometric pattern.
  - Solid lines pass through rows of atoms, representing planes in a crystalline structure.

- **Text and Labels:**
  - The label "Electron beam" accompanies an arrow pointing straight down towards the lattice, indicating the direction of an incident beam.
  - An angle labeled "θ" is shown between this incident beam and the line normal to atomic planes.
  - An arrow labeled "Scattered electrons" points away from the point of impact at an angle labeled "φ," representing the diffraction angle.
  - Distances are marked as "a" (between atoms in a row) and "d" (between atomic planes).

- **Relationships and Angles:**
  - The angles θ and φ demonstrate the beam angles concerning the atomic planes.
  - The incident and scattered beams are shown to indicate the interaction of electrons with the atomic lattice, resulting in diffraction.

This diagram effectively illustrates the principles of electron diffraction in a crystalline solid.

\#\# Dataset metadata

Quantum Phenomena; Physical Model Grounding Reasoning, Multi-Formula Reasoning

\paragraph{Recorded answer/context.}
D: D: \textbackslash{}( 0.218 \textbackslash{}text\{ nm\} \textbackslash{})

\paragraph{Figure/image path.}
/root/proposal\_for\_physic/science-mango/phyx\_mini\_run/phyx\_data/test\_image/532.png

\paragraph{Formalization target.}
create a compiling Lean file with sorry bodies at `PhyXMiniProblems/problem\_phyx\_mini\_0532.lean`. The Lean declarations must preserve the physical quantities, dimensions or dimensional roles, figure labels, governing-law hypotheses, and final relation expressed by this problem.
Use Mathlib/Physlib names found through LeanExplore where available. If a physics API is missing, introduce faithful local abstractions rather than scalar placeholder aliases.

\begin{theorem}[Physics formalization target]
\label{thm:physics:phyx_mini_0532:target}
\lean{PhyXMiniProblems.ProblemPhyXMini0532.problem_phyx_mini_0532}
\uses{def:physics:phyx-mini-0532:phyxminiproblems-problemphyxmini0532-lengthinnanometers, def:physics:phyx-mini-0532:phyxminiproblems-problemphyxmini0532-electrondiffractionsetup, def:physics:phyx-mini-0532:phyxminiproblems-problemphyxmini0532-matchesproblemreadouts, def:physics:phyx-mini-0532:phyxminiproblems-problemphyxmini0532-hasphysicalelectrondiffractionparameters, def:physics:phyx-mini-0532:phyxminiproblems-problemphyxmini0532-usesreferenceelectronconstants, def:physics:phyx-mini-0532:phyxminiproblems-problemphyxmini0532-matchessuppliedfigure, def:physics:phyx-mini-0532:phyxminiproblems-problemphyxmini0532-satisfiesdebroglieandbragglaws, def:physics:phyx-mini-0532:phyxminiproblems-problemphyxmini0532-isuniqueclosestdisplayedspacing}
The assigned autoformalize agent should translate this physics problem into Lean declarations in the covered file, with theorem and lemma proof bodies written as `by sorry`.
\end{theorem}
\begin{proof}
This is an autoformalization task, not a proof task. Produce statements that can later be proved by the physics prover without weakening the physical model.
\end{proof}

% --- Archon declaration topology begin ---
\section{Declaration topology}
\begin{definition}[Length Quantity]
  \label{def:physics:phyx-mini-0532:phyxminiproblems-problemphyxmini0532-lengthquantity}
  \lean{PhyXMiniProblems.ProblemPhyXMini0532.LengthQuantity}
  A nonnegative, unit-independent physical length
\end{definition}
\begin{proof}
  This is the declared model definition of Length Quantity.
\end{proof}
\begin{definition}[Mass Quantity]
  \label{def:physics:phyx-mini-0532:phyxminiproblems-problemphyxmini0532-massquantity}
  \lean{PhyXMiniProblems.ProblemPhyXMini0532.MassQuantity}
  A nonnegative, unit-independent physical mass
\end{definition}
\begin{proof}
  This is the declared model definition of Mass Quantity.
\end{proof}
\begin{definition}[action Dimension]
  \label{def:physics:phyx-mini-0532:phyxminiproblems-problemphyxmini0532-actiondimension}
  \lean{PhyXMiniProblems.ProblemPhyXMini0532.actionDimension}
  The dimension of action, used for the reduced Planck constant
\end{definition}
\begin{proof}
  This is the declared model definition of action Dimension.
\end{proof}
\begin{definition}[Action Quantity]
  \label{def:physics:phyx-mini-0532:phyxminiproblems-problemphyxmini0532-actionquantity}
  \lean{PhyXMiniProblems.ProblemPhyXMini0532.ActionQuantity}
  \uses{def:physics:phyx-mini-0532:phyxminiproblems-problemphyxmini0532-actiondimension}
  A nonnegative, unit-independent physical action
\end{definition}
\begin{proof}
  This is the declared model definition of Action Quantity.
\end{proof}
\begin{definition}[length Readout]
  \label{def:physics:phyx-mini-0532:phyxminiproblems-problemphyxmini0532-lengthreadout}
  \lean{PhyXMiniProblems.ProblemPhyXMini0532.lengthReadout}
  \uses{def:physics:phyx-mini-0532:phyxminiproblems-problemphyxmini0532-lengthquantity}
  Read a physical length in a selected Physlib length unit
\end{definition}
\begin{proof}
  This is the declared model definition of length Readout.
\end{proof}
\begin{definition}[length In Meters]
  \label{def:physics:phyx-mini-0532:phyxminiproblems-problemphyxmini0532-lengthinmeters}
  \lean{PhyXMiniProblems.ProblemPhyXMini0532.lengthInMeters}
  \uses{def:physics:phyx-mini-0532:phyxminiproblems-problemphyxmini0532-lengthquantity, def:physics:phyx-mini-0532:phyxminiproblems-problemphyxmini0532-lengthreadout}
  Metre readout of a physical length
\end{definition}
\begin{proof}
  This is the declared model definition of length In Meters.
\end{proof}
\begin{definition}[length In Nanometers]
  \label{def:physics:phyx-mini-0532:phyxminiproblems-problemphyxmini0532-lengthinnanometers}
  \lean{PhyXMiniProblems.ProblemPhyXMini0532.lengthInNanometers}
  \uses{def:physics:phyx-mini-0532:phyxminiproblems-problemphyxmini0532-lengthquantity, def:physics:phyx-mini-0532:phyxminiproblems-problemphyxmini0532-lengthreadout}
  Nanometre readout of a physical length
\end{definition}
\begin{proof}
  This is the declared model definition of length In Nanometers.
\end{proof}
\begin{definition}[energy In Joules]
  \label{def:physics:phyx-mini-0532:phyxminiproblems-problemphyxmini0532-energyinjoules}
  \lean{PhyXMiniProblems.ProblemPhyXMini0532.energyInJoules}
  Coherent-SI readout of an energy, in joules
\end{definition}
\begin{proof}
  This is the declared model definition of energy In Joules.
\end{proof}
\begin{definition}[electron Volt In Joules]
  \label{def:physics:phyx-mini-0532:phyxminiproblems-problemphyxmini0532-electronvoltinjoules}
  \lean{PhyXMiniProblems.ProblemPhyXMini0532.electronVoltInJoules}
  \uses{def:physics:phyx-mini-0532:phyxminiproblems-problemphyxmini0532-energyinjoules}
  The joule readout of Physlib's dimensionful electron volt
\end{definition}
\begin{proof}
  This is the declared model definition of electron Volt In Joules.
\end{proof}
\begin{definition}[energy In Electron Volts]
  \label{def:physics:phyx-mini-0532:phyxminiproblems-problemphyxmini0532-energyinelectronvolts}
  \lean{PhyXMiniProblems.ProblemPhyXMini0532.energyInElectronVolts}
  \uses{def:physics:phyx-mini-0532:phyxminiproblems-problemphyxmini0532-energyinjoules, def:physics:phyx-mini-0532:phyxminiproblems-problemphyxmini0532-electronvoltinjoules}
  Read a dimensionful energy as a scalar number of electron volts
\end{definition}
\begin{proof}
  This is the declared model definition of energy In Electron Volts.
\end{proof}
\begin{definition}[mass In Kilograms]
  \label{def:physics:phyx-mini-0532:phyxminiproblems-problemphyxmini0532-massinkilograms}
  \lean{PhyXMiniProblems.ProblemPhyXMini0532.massInKilograms}
  \uses{def:physics:phyx-mini-0532:phyxminiproblems-problemphyxmini0532-massquantity}
  Coherent-SI readout of a mass, in kilograms
\end{definition}
\begin{proof}
  This is the declared model definition of mass In Kilograms.
\end{proof}
\begin{definition}[action In Joule Seconds]
  \label{def:physics:phyx-mini-0532:phyxminiproblems-problemphyxmini0532-actioninjouleseconds}
  \lean{PhyXMiniProblems.ProblemPhyXMini0532.actionInJouleSeconds}
  \uses{def:physics:phyx-mini-0532:phyxminiproblems-problemphyxmini0532-actionquantity}
  Coherent-SI readout of an action, in joule-seconds
\end{definition}
\begin{proof}
  This is the declared model definition of action In Joule Seconds.
\end{proof}
\begin{definition}[Beam Particle]
  \label{def:physics:phyx-mini-0532:phyxminiproblems-problemphyxmini0532-beamparticle}
  \lean{PhyXMiniProblems.ProblemPhyXMini0532.BeamParticle}
  Species of particle in the incident beam
\end{definition}
\begin{proof}
  This is the declared model definition of Beam Particle.
\end{proof}
\begin{definition}[Crystal Material]
  \label{def:physics:phyx-mini-0532:phyxminiproblems-problemphyxmini0532-crystalmaterial}
  \lean{PhyXMiniProblems.ProblemPhyXMini0532.CrystalMaterial}
  Crystal material used as the diffraction target
\end{definition}
\begin{proof}
  This is the declared model definition of Crystal Material.
\end{proof}
\begin{definition}[Figure Label]
  \label{def:physics:phyx-mini-0532:phyxminiproblems-problemphyxmini0532-figurelabel}
  \lean{PhyXMiniProblems.ProblemPhyXMini0532.FigureLabel}
  Text or symbolic labels visible in the supplied lattice diagram
\end{definition}
\begin{proof}
  This is the declared model definition of Figure Label.
\end{proof}
\begin{definition}[Figure Object]
  \label{def:physics:phyx-mini-0532:phyxminiproblems-problemphyxmini0532-figureobject}
  \lean{PhyXMiniProblems.ProblemPhyXMini0532.FigureObject}
  Geometrical objects drawn in the supplied lattice diagram
\end{definition}
\begin{proof}
  This is the declared model definition of Figure Object.
\end{proof}
\begin{definition}[Davisson Germer Figure]
  \label{def:physics:phyx-mini-0532:phyxminiproblems-problemphyxmini0532-davissongermerfigure}
  \lean{PhyXMiniProblems.ProblemPhyXMini0532.DavissonGermerFigure}
  \uses{def:physics:phyx-mini-0532:phyxminiproblems-problemphyxmini0532-figurelabel, def:physics:phyx-mini-0532:phyxminiproblems-problemphyxmini0532-figureobject}
  Presentation-level data from the primary figure. It records only which labels and objects are shown; their physical quantities live in the experiment setup
\end{definition}
\begin{proof}
  This is the declared model definition of Davisson Germer Figure.
\end{proof}
\begin{definition}[Electron Diffraction Setup]
  \label{def:physics:phyx-mini-0532:phyxminiproblems-problemphyxmini0532-electrondiffractionsetup}
  \lean{PhyXMiniProblems.ProblemPhyXMini0532.ElectronDiffractionSetup}
  \uses{def:physics:phyx-mini-0532:phyxminiproblems-problemphyxmini0532-lengthquantity, def:physics:phyx-mini-0532:phyxminiproblems-problemphyxmini0532-massquantity, def:physics:phyx-mini-0532:phyxminiproblems-problemphyxmini0532-actionquantity, def:physics:phyx-mini-0532:phyxminiproblems-problemphyxmini0532-beamparticle, def:physics:phyx-mini-0532:phyxminiproblems-problemphyxmini0532-crystalmaterial, def:physics:phyx-mini-0532:phyxminiproblems-problemphyxmini0532-davissongermerfigure}
  The physical experiment and all quantities named in the prose or figure. Angles use Real.Angle; toReal supplies the principal radian representative needed for the pictured acute-angle geometry
\end{definition}
\begin{proof}
  This is the declared model definition of Electron Diffraction Setup.
\end{proof}
\begin{definition}[Matches Problem Readouts]
  \label{def:physics:phyx-mini-0532:phyxminiproblems-problemphyxmini0532-matchesproblemreadouts}
  \lean{PhyXMiniProblems.ProblemPhyXMini0532.MatchesProblemReadouts}
  \uses{def:physics:phyx-mini-0532:phyxminiproblems-problemphyxmini0532-energyinelectronvolts, def:physics:phyx-mini-0532:phyxminiproblems-problemphyxmini0532-electrondiffractionsetup}
  Numerical and categorical data stated in the prose: electrons strike nickel with kinetic energy 54.0 eV, and the observed first maximum has phi = 50 degrees = 5*pi/18 radians. No value of a is assumed
\end{definition}
\begin{proof}
  This is the declared model definition of Matches Problem Readouts.
\end{proof}
\begin{definition}[Has Physical Electron Diffraction Parameters]
  \label{def:physics:phyx-mini-0532:phyxminiproblems-problemphyxmini0532-hasphysicalelectrondiffractionparameters}
  \lean{PhyXMiniProblems.ProblemPhyXMini0532.HasPhysicalElectronDiffractionParameters}
  \uses{def:physics:phyx-mini-0532:phyxminiproblems-problemphyxmini0532-lengthinmeters, def:physics:phyx-mini-0532:phyxminiproblems-problemphyxmini0532-energyinjoules, def:physics:phyx-mini-0532:phyxminiproblems-problemphyxmini0532-massinkilograms, def:physics:phyx-mini-0532:phyxminiproblems-problemphyxmini0532-actioninjouleseconds, def:physics:phyx-mini-0532:phyxminiproblems-problemphyxmini0532-electrondiffractionsetup}
  Physical sign and branch conditions. The angular bounds select the pictured acute incidence angle and non-reflex scattering angle without determining the requested lattice spacing
\end{definition}
\begin{proof}
  This is the declared model definition of Has Physical Electron Diffraction Parameters.
\end{proof}
\begin{definition}[Uses Reference Electron Constants]
  \label{def:physics:phyx-mini-0532:phyxminiproblems-problemphyxmini0532-usesreferenceelectronconstants}
  \lean{PhyXMiniProblems.ProblemPhyXMini0532.UsesReferenceElectronConstants}
  \uses{def:physics:phyx-mini-0532:phyxminiproblems-problemphyxmini0532-massinkilograms, def:physics:phyx-mini-0532:phyxminiproblems-problemphyxmini0532-actioninjouleseconds, def:physics:phyx-mini-0532:phyxminiproblems-problemphyxmini0532-electrondiffractionsetup}
  Calibration of the electron constants used by the de Broglie relation. Physlib supplies the reduced Planck constant in joule-seconds. Since no electron-rest-mass declaration was found, its kilogram readout is explicit
\end{definition}
\begin{proof}
  This is the declared model definition of Uses Reference Electron Constants.
\end{proof}
\begin{definition}[Matches Supplied Figure]
  \label{def:physics:phyx-mini-0532:phyxminiproblems-problemphyxmini0532-matchessuppliedfigure}
  \lean{PhyXMiniProblems.ProblemPhyXMini0532.MatchesSuppliedFigure}
  \uses{def:physics:phyx-mini-0532:phyxminiproblems-problemphyxmini0532-lengthreadout, def:physics:phyx-mini-0532:phyxminiproblems-problemphyxmini0532-figurelabel, def:physics:phyx-mini-0532:phyxminiproblems-problemphyxmini0532-figureobject, def:physics:phyx-mini-0532:phyxminiproblems-problemphyxmini0532-electrondiffractionsetup}
  The qualitative labels and objects, together with the quantitative relations read from the diagram. Since phi is measured from the backward extension of the incident beam, 2*theta + phi = pi. Since a is horizontal while d is perpendicular to the slanted planes, d = a*cos(theta) in every length unit
\end{definition}
\begin{proof}
  This is the declared model definition of Matches Supplied Figure.
\end{proof}
\begin{definition}[Satisfies De Broglie And Bragg Laws]
  \label{def:physics:phyx-mini-0532:phyxminiproblems-problemphyxmini0532-satisfiesdebroglieandbragglaws}
  \lean{PhyXMiniProblems.ProblemPhyXMini0532.SatisfiesDeBroglieAndBraggLaws}
  \uses{def:physics:phyx-mini-0532:phyxminiproblems-problemphyxmini0532-lengthreadout, def:physics:phyx-mini-0532:phyxminiproblems-problemphyxmini0532-lengthinmeters, def:physics:phyx-mini-0532:phyxminiproblems-problemphyxmini0532-energyinjoules, def:physics:phyx-mini-0532:phyxminiproblems-problemphyxmini0532-massinkilograms, def:physics:phyx-mini-0532:phyxminiproblems-problemphyxmini0532-actioninjouleseconds, def:physics:phyx-mini-0532:phyxminiproblems-problemphyxmini0532-electrondiffractionsetup}
  The electron wave obeys the nonrelativistic de Broglie relation lambda = 2*pi*hbar/sqrt(2*m*K). Constructive interference from adjacent atomic planes obeys Bragg's law n*lambda = 2*d*sin(theta). These are governing laws and contain no numerical lattice-spacing answer
\end{definition}
\begin{proof}
  This is the declared model definition of Satisfies De Broglie And Bragg Laws.
\end{proof}
\begin{definition}[Answer Choice]
  \label{def:physics:phyx-mini-0532:phyxminiproblems-problemphyxmini0532-answerchoice}
  \lean{PhyXMiniProblems.ProblemPhyXMini0532.AnswerChoice}
  Labels of the four answer choices supplied with the problem
\end{definition}
\begin{proof}
  This is the declared model definition of Answer Choice.
\end{proof}
\begin{definition}[spacing In Nanometers]
  \label{def:physics:phyx-mini-0532:phyxminiproblems-problemphyxmini0532-answerchoice-spacinginnanometers}
  \lean{PhyXMiniProblems.ProblemPhyXMini0532.AnswerChoice.spacingInNanometers}
  \uses{def:physics:phyx-mini-0532:phyxminiproblems-problemphyxmini0532-answerchoice}
  Lattice-spacing value printed beside each answer, in nanometres
\end{definition}
\begin{proof}
  This is the declared model definition of spacing In Nanometers.
\end{proof}
\begin{definition}[recorded Dataset Answer]
  \label{def:physics:phyx-mini-0532:phyxminiproblems-problemphyxmini0532-recordeddatasetanswer}
  \lean{PhyXMiniProblems.ProblemPhyXMini0532.recordedDatasetAnswer}
  \uses{def:physics:phyx-mini-0532:phyxminiproblems-problemphyxmini0532-answerchoice}
  The answer label recorded by the source dataset
\end{definition}
\begin{proof}
  This is the declared model definition of recorded Dataset Answer.
\end{proof}
\begin{definition}[Is Unique Closest Displayed Spacing]
  \label{def:physics:phyx-mini-0532:phyxminiproblems-problemphyxmini0532-isuniqueclosestdisplayedspacing}
  \lean{PhyXMiniProblems.ProblemPhyXMini0532.IsUniqueClosestDisplayedSpacing}
  \uses{def:physics:phyx-mini-0532:phyxminiproblems-problemphyxmini0532-lengthinnanometers, def:physics:phyx-mini-0532:phyxminiproblems-problemphyxmini0532-electrondiffractionsetup, def:physics:phyx-mini-0532:phyxminiproblems-problemphyxmini0532-answerchoice}
  A choice is uniquely closest to the lattice spacing inferred from the model. This compares the inferred value with source metadata; it is not an assumed experimental law
\end{definition}
\begin{proof}
  This is the declared model definition of Is Unique Closest Displayed Spacing.
\end{proof}
% --- Archon declaration topology end ---
% --- Archon physics formalization source end ---
